% ============================================================
%  Nota di supporto a dispense/backup/02c_eiopa_rfr_smith_wilson.tex
%  OBIETTIVO: dimostrare che la matrice di Gram H = [H(u_i,u_j)] del
%  problema Smith-Wilson e' simmetrica definita positiva (SPD).
%  Dimostrazione centrata sul Teorema di Mercer.
%  Documento standalone, compilabile con pdflatex.
% ============================================================
\documentclass[11pt,a4paper]{article}
\usepackage[T1]{fontenc}
\usepackage[utf8]{inputenc}
\usepackage[italian]{babel}
\usepackage{amsmath,amssymb,amsthm,mathtools}
\usepackage{geometry}
\geometry{margin=2.7cm}
\usepackage{xcolor}
\usepackage{hyperref}
\hypersetup{colorlinks=true,linkcolor=blue!55!black,citecolor=blue!55!black,urlcolor=blue!55!black}

\theoremstyle{plain}
\newtheorem{teorema}{Teorema}
\newtheorem{proposizione}[teorema]{Proposizione}
\newtheorem{lemma}[teorema]{Lemma}
\newtheorem{corollario}[teorema]{Corollario}
\theoremstyle{remark}
\newtheorem*{osservazione}{Osservazione}

\DeclareMathOperator{\tr}{tr}
\newcommand{\Hmat}{\mathbf{H}}
\newcommand{\Qmat}{\mathbf{Q}}
\newcommand{\transp}{^{\!\top}}
\newcommand{\LLP}{\mathrm{LLP}}
\newcommand{\R}{\mathbb{R}}
\newcommand{\Top}{T_{\!H}}
\newcommand{\inner}[2]{\langle #1, #2\rangle}

\title{\textbf{La matrice di Gram del problema Smith--Wilson\\[2pt]
\`e definita positiva}\\[6pt]
\large Nota di supporto a \texttt{dispense/backup/02c\_eiopa\_rfr\_smith\_wilson.tex} --- dimostrazione via il
Teorema di Mercer}
\author{}
\date{}

\begin{document}
\maketitle

% ============================================================
\section{Obiettivo e mappa della dimostrazione}
\label{sec:obiettivo}
% ============================================================

La dispensa principale costruisce la funzione di sconto Smith--Wilson
$P(t)=e^{-\omega t}\bigl(1+\sum_{j=1}^m \zeta_j\,H(t,u_j)\bigr)$ risolvendo un \textbf{problema
di minimo vincolato} nello spazio finito-dimensionale dei coefficienti $\zeta\in\R^m$:
\begin{equation}
  \min_{\zeta\in\R^m}\ \zeta\transp\Hmat\,\zeta
  \qquad\text{soggetto ai vincoli di pricing}\qquad
  \Qmat\transp\Hmat\,\zeta = p-q,
  \label{eq:min}
\end{equation}
dove $\Hmat=[H(u_i,u_j)]_{i,j=1}^m\in\R^{m\times m}$ \`e la \textbf{matrice di Gram} del
\emph{nucleo di Wilson}
\begin{equation}
  H(t,u) = \alpha\,\min(t,u) - e^{-\alpha\,\max(t,u)}\,\sinh\!\bigl(\alpha\,\min(t,u)\bigr),
  \qquad \alpha>0,\ t,u\ge0,
  \label{eq:H-def}
\end{equation}
valutata sugli $m$ nodi di pagamento distinti $u_1<\cdots<u_m$ (nella dispensa $u_j=j$,
$m=\LLP=20$).

L'intera costruzione dipende da un'unica propriet\`a algebrica, che la dispensa \emph{assume}
senza dimostrare (per non usare l'analisi funzionale in un corso triennale). Questa nota la
dimostra.

\begin{osservazione}
Questa \`e la \textbf{terza} via seguita nel materiale del corso per la stessa propriet\`a: la
nota \texttt{dispense/backup/02e\_eiopa\_rfr\_smith\_wilson.tex} la ricava per via
\emph{variazionale} (il nucleo di Wilson come rappresentante di Riesz di un prodotto scalare,
quindi $\Hmat$ come matrice di Gram in senso proprio), mentre le note
\texttt{dispense/backup/02f\_eiopa\_rfr\_smith\_wilson.tex} e
\texttt{dispense/backup/02g\_eiopa\_rfr\_smith\_wilson.tex} la dimostrano in modo
\emph{elementare} (rappresentazione integrale del nucleo, induzione discendente sui nodi, senza
analisi funzionale). Qui si segue invece il \textbf{Teorema di Mercer}.
\end{osservazione}

\begin{proposizione}[Perch\'e serve la definita positivit\`a]
\label{prop:benposto}
Se $\Hmat$ \`e simmetrica \textbf{definita positiva} (SPD), allora $\zeta\mapsto\zeta\transp\Hmat\zeta$
\`e strettamente convessa: il problema~\eqref{eq:min} ha \textbf{un unico} minimo, il sistema di
calibrazione $(\Qmat\transp\Hmat\Qmat)\,b=p-q$ ha \textbf{soluzione unica} ed \`e risolubile in
modo stabile con \textbf{Cholesky}. Se invece $\Hmat$ avesse un autovalore $\le0$, il minimo
potrebbe non esistere o non essere unico e Cholesky fallirebbe.
\end{proposizione}

L'obiettivo della nota \`e dunque il seguente teorema.

\begin{teorema}[Obiettivo]
\label{thm:obiettivo}
Per ogni scelta di nodi \textbf{distinti} $u_1<\cdots<u_m$ con $u_i>0$, la matrice di Gram
$\Hmat=[H(u_i,u_j)]_{i,j=1}^m$ \`e simmetrica definita positiva.
\end{teorema}

\paragraph{Mappa della dimostrazione (una sola catena lineare).}
\begin{center}
\fbox{\parbox{0.92\textwidth}{\centering
rappresentazione spettrale di $H$ \textbf{(§\ref{sec:fourier})}
$\;\Rightarrow\;$ l'operatore integrale $\Top$ \`e positivo \textbf{(§\ref{sec:operatore})}
$\;\Rightarrow\;$ espansione di Mercer $H=\sum_n\lambda_n\phi_n\phi_n$ \textbf{(§\ref{sec:mercer})}
$\;\Rightarrow\;$ $\Hmat$ \`e SPD \textbf{(§\ref{sec:spd})}.}}
\end{center}
Il \textbf{Teorema di Mercer} (§\ref{sec:mercer}) \`e il cuore della dimostrazione. L'unico
strumento analitico \`e la rappresentazione spettrale di $H$ (§\ref{sec:fourier}), usata in
\emph{due} punti chiaramente segnalati: per verificare l'ipotesi di Mercer (§\ref{sec:operatore})
e per la stretta positivit\`a (§\ref{sec:spd}).

% ============================================================
\section{Lo strumento: la rappresentazione spettrale di \texorpdfstring{$H$}{H}}
\label{sec:fourier}
% ============================================================

\begin{lemma}[Riscrittura algebrica]
\label{lem:algebrico}
Per ogni $t,u\ge0$:\quad
$H(t,u) = \alpha\min(t,u) - \tfrac12 e^{-\alpha|t-u|} + \tfrac12 e^{-\alpha(t+u)}.$
\end{lemma}
\begin{proof}
Con $m=\min(t,u)$, $M=\max(t,u)$: $M-m=|t-u|$, $M+m=t+u$, e da
$\sinh(\alpha m)=\tfrac12(e^{\alpha m}-e^{-\alpha m})$ segue
$e^{-\alpha M}\sinh(\alpha m)=\tfrac12(e^{-\alpha|t-u|}-e^{-\alpha(t+u)})$; sostituire
in~\eqref{eq:H-def}.
\end{proof}

Servono due identit\`a integrali classiche di Fourier (tavole standard, es.\
\cite{gradshteyn2007,bracewell2000}):
\begin{equation}
  \frac{2}{\pi}\!\int_0^\infty\! \frac{\sin\omega t\,\sin\omega u}{\omega^2}\,d\omega = \min(t,u),
  \qquad
  \frac{2}{\pi}\!\int_0^\infty\! \frac{\cos\omega x}{\omega^2+\alpha^2}\,d\omega = \frac{1}{\alpha}e^{-\alpha|x|}.
  \label{eq:fourier}
\end{equation}

\begin{lemma}[Rappresentazione spettrale]
\label{lem:spettrale}
\[
  H(t,u) = \int_0^\infty \rho(\omega)\,\sin(\omega t)\sin(\omega u)\,d\omega,
  \qquad
  \rho(\omega):=\frac{2\alpha^3}{\pi\,\omega^2(\omega^2+\alpha^2)}\;>\;0\quad\forall\,\omega>0.
\]
\end{lemma}
\begin{proof}
Per fratti semplici $\tfrac{1}{\omega^2(\omega^2+\alpha^2)}=\tfrac1{\alpha^2}\bigl(\tfrac1{\omega^2}
-\tfrac1{\omega^2+\alpha^2}\bigr)$. Con $\sin\omega t\,\sin\omega u=\tfrac12[\cos\omega(t-u)-\cos\omega(t+u)]$
e la seconda identit\`a~\eqref{eq:fourier} (per $x=t-u$ e $x=t+u\ge0$),
$\int_0^\infty\tfrac{\sin\omega t\,\sin\omega u}{\omega^2+\alpha^2}d\omega=\tfrac{\pi}{4\alpha}
(e^{-\alpha|t-u|}-e^{-\alpha(t+u)})$; con la prima identit\`a per l'altro addendo e moltiplicando
per $\tfrac{2\alpha^3}{\pi}$, il confronto con il Lemma~\ref{lem:algebrico} d\`a la tesi.
\end{proof}

\begin{osservazione}[Il perno di tutta la nota]
La densit\`a spettrale $\rho$ \`e \textbf{ovunque positiva}: sebbene $H$ contenga un segno meno
(Lemma~\ref{lem:algebrico}), scritto come integrale con peso $\rho\ge0$ non c'\`e cancellazione.
\`E questa positivit\`a l'unico ingrediente analitico della dimostrazione, usato in
§\ref{sec:operatore} e in §\ref{sec:spd}.
\end{osservazione}

% ============================================================
\section{L'operatore integrale \texorpdfstring{$\Top$}{T\_H} \`e positivo}
\label{sec:operatore}
% ============================================================

Sia $L:=\LLP$ e $X:=[0,L]$, compatto contenente tutti i nodi. Associamo a $H$ l'operatore
integrale su $L^2(X)$
\[
  \Top f(s):=\int_0^L H(s,t)\,f(t)\,dt .
\]
Poich\'e $H$ \`e \textbf{continuo} su $X\times X$ (le due espressioni per $t\lessgtr u$ coincidono
in $t=u$) e \textbf{simmetrico}, l'operatore $\Top$ \`e di Hilbert--Schmidt, dunque
\textbf{compatto}, ed \textbf{autoaggiunto} \cite{reedsimon1980,brezis2011}. Resta la propriet\`a
non banale, che \`e proprio l'ipotesi del Teorema di Mercer.

\begin{proposizione}[Positivit\`a dell'operatore]
\label{prop:positivo}
$\Top$ \`e un \textbf{operatore positivo}: $\inner{\Top f}{f}_{L^2}\ge0$ per ogni $f\in L^2(X)$.
\end{proposizione}
\begin{proof}
Inserendo la rappresentazione spettrale (Lemma~\ref{lem:spettrale}) e scambiando l'ordine di
integrazione (Fubini--Tonelli, lecito perch\'e $H$ \`e continuo sul compatto $X\times X$):
\[
  \inner{\Top f}{f}_{L^2}
  = \int_0^L\!\!\int_0^L H(s,t)f(s)f(t)\,ds\,dt
  = \int_0^\infty \rho(\omega)\Bigl(\int_0^L f(t)\sin(\omega t)\,dt\Bigr)^{\!2} d\omega\ \ge\ 0,
\]
essendo $\rho(\omega)\ge0$ e il fattore restante un quadrato.
\end{proof}

% ============================================================
\section{Il Teorema di Mercer}
\label{sec:mercer}
% ============================================================

Essendo $\Top$ compatto e autoaggiunto, per il \textbf{teorema spettrale per operatori compatti
autoaggiunti} \cite{reedsimon1980,brezis2011} esiste una base ortonormale $\{\phi_n\}_{n\ge1}$ di
$L^2(X)$ di autofunzioni, $\Top\phi_n=\lambda_n\phi_n$, con autovalori reali $\lambda_n\to0$. Per
la Proposizione~\ref{prop:positivo} gli autovalori sono \textbf{non negativi},
$\lambda_n=\inner{\Top\phi_n}{\phi_n}\ge0$; inoltre ogni autofunzione con $\lambda_n\ne0$ \`e
\textbf{continua} ($\phi_n=\lambda_n^{-1}\Top\phi_n$ e $\Top$ manda $L^2(X)$ in $C(X)$).

\begin{teorema}[Mercer, 1909 \cite{mercer1909}; cfr.\ \cite{riesznagy1955,steinwart2008}]
\label{thm:mercer}
Sia $K$ continuo e simmetrico sul compatto $X\times X$, con operatore integrale $T_K$ positivo su
$L^2(X)$. Allora
\[
  K(s,t)=\sum_{n=1}^\infty \lambda_n\,\phi_n(s)\,\phi_n(t),
\]
con $\lambda_n\ge0$, $\{\phi_n\}$ autofunzioni continue ortonormali, e convergenza
\textbf{assoluta e uniforme} su $X\times X$.
\end{teorema}

\begin{osservazione}[Ingredienti]
La dimostrazione (che non riportiamo, cfr.\ \cite{riesznagy1955,steinwart2008}) usa la positivit\`a
dell'operatore per rendere i resti $\ge0$ e decrescenti sulla diagonale; il \textbf{Teorema di
Dini} d\`a l'uniformit\`a sulla diagonale, Cauchy--Schwarz fuori diagonale. Ne segue la
\textbf{traccia} $\sum_n\lambda_n=\int_0^L H(t,t)\,dt=\tr\Top<\infty$ (finita perch\'e $H$ \`e
continuo sul compatto), quindi gli autovalori sono sommabili.
\end{osservazione}

Il nucleo $H$ soddisfa tutte le ipotesi (continuo, simmetrico, $\Top$ positivo per la
Proposizione~\ref{prop:positivo}), dunque ammette l'espansione di Mercer
$H(s,t)=\sum_n\lambda_n\phi_n(s)\phi_n(t)$.

% ============================================================
\section{La matrice di Gram \texorpdfstring{$\Hmat$}{H} \`e SPD}
\label{sec:spd}
% ============================================================

Concludiamo la dimostrazione del Teorema~\ref{thm:obiettivo}. La strada \`e diretta e usa
\emph{un fatto per volta}:
\begin{itemize}
  \item la \textbf{simmetria} \`e ovvia da $H(t,u)=H(u,t)$;
  \item la \textbf{semidefinita} positivit\`a ($\zeta\transp\Hmat\zeta\ge0$) segue dall'espansione
    di Mercer (§\ref{sec:mercer});
  \item la \textbf{stretta} positivit\`a ($\zeta\transp\Hmat\zeta>0$ per $\zeta\ne0$) segue dalla
    rappresentazione di Fourier (§\ref{sec:fourier}) e dal Lemma~\ref{lem:indip} qui sotto.
\end{itemize}
Isoliamo prima il fatto di algebra lineare che serve alla stretta positivit\`a.

\begin{lemma}[Indipendenza lineare dei seni campionati]
\label{lem:indip}
Siano $u_1<\cdots<u_m$ nodi \emph{distinti} con $u_i>0$ e sia $c\in\R^m$. Se
\[
  g(\omega):=\sum_{i=1}^m c_i\,\sin(\omega u_i)=0 \qquad\text{per ogni }\omega\ge0,
\]
allora $c=0$.
\end{lemma}
\begin{proof}
La funzione $g$ \`e identicamente nulla, quindi tutte le sue derivate in $\omega=0$ sono nulle.
Poich\'e $\dfrac{d^{2k+1}}{d\omega^{2k+1}}\sin(\omega u_i)\Big|_{\omega=0}=(-1)^k u_i^{2k+1}$
(le derivate di ordine pari si annullano in $0$), la derivata di ordine $2k+1$ di $g$ in $0$ d\`a
\[
  \sum_{i=1}^m c_i\,u_i^{2k+1}=0 \qquad\text{per ogni }k=0,1,2,\dots
\]
Poniamo $y_i:=c_i u_i$ e prendiamo le prime $m$ equazioni ($k=0,\dots,m-1$):
\[
  \sum_{i=1}^m y_i\,(u_i^2)^{k}=0,\qquad k=0,\dots,m-1,
\]
cio\`e il sistema lineare $V y=0$, dove $V=\bigl[(u_i^2)^{k}\bigr]$ \`e la matrice di
\textbf{Vandermonde} costruita sui valori $u_1^2,\dots,u_m^2$. Essendo gli $u_i>0$ distinti,
anche gli $u_i^2$ sono distinti, quindi $V$ \`e invertibile \cite{quarteroni2017}: $y=0$. Infine
$c_i=y_i/u_i=0$ per ogni $i$.
\end{proof}

\begin{proof}[Dimostrazione del Teorema~\ref{thm:obiettivo}]
\textbf{Simmetria.} $\Hmat_{ij}=H(u_i,u_j)=H(u_j,u_i)=\Hmat_{ji}$.

\textbf{Semidefinita positivit\`a.} Sia $\zeta\in\R^m$. Sostituiamo nell'espressione di
$\zeta\transp\Hmat\zeta$ l'espansione di Mercer $H(u_i,u_j)=\sum_n\lambda_n\phi_n(u_i)\phi_n(u_j)$
(§\ref{sec:mercer}); lo scambio tra la somma \emph{finita} su $i,j$ e la serie \`e lecito per la
convergenza uniforme:
\[
  \zeta\transp\Hmat\zeta
  = \sum_{i,j=1}^m \zeta_i\zeta_j\,H(u_i,u_j)
  = \sum_{n=1}^\infty \lambda_n\Bigl(\underbrace{\sum_{j=1}^m \zeta_j\,\phi_n(u_j)}_{=:v_n}\Bigr)^{\!2}
  = \sum_{n=1}^\infty \lambda_n\,v_n^2\ \ge\ 0,
\]
poich\'e ogni $\lambda_n\ge0$ e ogni $v_n^2\ge0$.

\textbf{Stretta positivit\`a.} Supponiamo $\zeta\transp\Hmat\zeta=0$ e mostriamo $\zeta=0$.
Conviene ripartire dalla rappresentazione di Fourier (Lemma~\ref{lem:spettrale}): sostituendo
$H(u_i,u_j)=\int_0^\infty\rho(\omega)\sin(\omega u_i)\sin(\omega u_j)\,d\omega$ e portando fuori
la somma finita,
\[
  0=\zeta\transp\Hmat\zeta
   =\int_0^\infty \rho(\omega)\Bigl(\sum_{i=1}^m \zeta_i\sin(\omega u_i)\Bigr)^{\!2} d\omega .
\]
L'integrando \`e continuo e non negativo (perch\'e $\rho>0$); un integrale nullo di una funzione
continua non negativa impone che l'integrando sia identicamente nullo, dunque
$\sum_{i=1}^m \zeta_i\sin(\omega u_i)=0$ per ogni $\omega\ge0$. Per il Lemma~\ref{lem:indip},
$\zeta=0$.

I tre punti danno che $\Hmat$ \`e simmetrica definita positiva.
\end{proof}

\begin{osservazione}[Lettura spettrale]
La parte semidefinita mostra $\Hmat=\sum_{n=1}^\infty\lambda_n\,v_nv_n\transp$ con
$v_n=(\phi_n(u_1),\dots,\phi_n(u_m))\transp\in\R^m$: la matrice di Gram \`e il
\textbf{campionamento sui nodi} della decomposizione spettrale del nucleo continuo, l'analogo
infinito-dimensionale della decomposizione $\mathbf A=\sum_k\lambda_k v_kv_k\transp$ di una
matrice simmetrica.
\end{osservazione}

\begin{corollario}
Per la Proposizione~\ref{prop:benposto}, il problema di minimo~\eqref{eq:min} ha soluzione unica,
il sistema di calibrazione $(\Qmat\transp\Hmat\Qmat)\,b=p-q$ \`e univocamente risolubile e la sua
soluzione via Cholesky \`e numericamente stabile.
\end{corollario}

\begin{osservazione}[Il sistema completo]
Il sistema effettivamente risolto ha matrice $\Qmat\transp\Hmat\Qmat$, SPD non solo perch\'e
$\Hmat$ lo \`e (questa nota) ma anche perch\'e $\Qmat$ ha rango pieno di colonna: quest'ultimo
fatto, di sola algebra lineare, \`e dimostrato nella nota gemella \texttt{note\_QHQ\_spd.txt}.
\end{osservazione}

% ============================================================
\section*{Nota sul livello di rigore}
% ============================================================
\addcontentsline{toc}{section}{Nota sul livello di rigore}
La riscrittura algebrica (Lemma~\ref{lem:algebrico}) e la rappresentazione spettrale
(Lemma~\ref{lem:spettrale}) sono state verificate, in fase di stesura, sia analiticamente sia
numericamente. Sono citati senza ridimostrazione i risultati standard: le identit\`a di
Fourier~\eqref{eq:fourier}, il teorema spettrale per operatori compatti autoaggiunti, il Teorema
di Dini e il Teorema di Mercer. La strategia (nucleo definito positivo / decomposizione spettrale)
\`e quella classica per questa classe di kernel \cite{aronszajn1950,wahba1990,fasshauer2007};
non ho verificato i numeri di pagina/teorema delle fonti citate.

% ============================================================
\begin{thebibliography}{99}

\bibitem{aronszajn1950}
\textbf{Aronszajn, N.} (1950). Theory of Reproducing Kernels.
\textit{Transactions of the American Mathematical Society}, 68(3), 337--404.

\bibitem{brezis2011}
\textbf{Brezis, H.} (2011). \textit{Functional Analysis, Sobolev Spaces and Partial Differential
Equations}. Springer. (Operatori compatti e teorema spettrale: cap.~6.)

\bibitem{bracewell2000}
\textbf{Bracewell, R.N.} (2000). \textit{The Fourier Transform and Its Applications}, 3rd ed.
McGraw-Hill.

\bibitem{fasshauer2007}
\textbf{Fasshauer, G.E.} (2007). \textit{Meshfree Approximation Methods with MATLAB}. World
Scientific.

\bibitem{gradshteyn2007}
\textbf{Gradshteyn, I.S., Ryzhik, I.M.} (2007). \textit{Table of Integrals, Series, and Products},
7th ed. Academic Press.

\bibitem{mercer1909}
\textbf{Mercer, J.} (1909). Functions of positive and negative type, and their connection with the
theory of integral equations. \textit{Philosophical Transactions of the Royal Society A}, 209,
415--446.

\bibitem{quarteroni2017}
\textbf{Quarteroni, A., Saleri, F., Gervasio, P.} (2017). \textit{Calcolo Scientifico}, 6a ed.
Springer.

\bibitem{reedsimon1980}
\textbf{Reed, M., Simon, B.} (1980). \textit{Methods of Modern Mathematical Physics I: Functional
Analysis}, rev.\ ed. Academic Press.

\bibitem{riesznagy1955}
\textbf{Riesz, F., Sz.-Nagy, B.} (1955). \textit{Functional Analysis}. Frederick Ungar.

\bibitem{steinwart2008}
\textbf{Steinwart, I., Christmann, A.} (2008). \textit{Support Vector Machines}. Springer.
(Teorema di Mercer: \S4.)

\bibitem{wahba1990}
\textbf{Wahba, G.} (1990). \textit{Spline Models for Observational Data}. CBMS-NSF Regional
Conference Series in Applied Mathematics, Vol.~59. SIAM.

\end{thebibliography}

\end{document}
